\chapter{Hints}

\section{Mechanics}

\begin{hint}{1.01}
First find the initial speed $v_0$ required for the ball to reach the top of the wall as a function of the angle of projection $\theta$; then find the minimum of $v_0$.
\end{hint}

\begin{hint}{1.02}
Treat the motion of the ball above the tops of the walls as an independent projectile motion; then relate the motion above the wall tops to the motion below them.
\end{hint}

\begin{hint}{1.03}
Discuss the condition for the trajectory of the shell to be tangent to the shelter.
\end{hint}

\begin{hint}{1.04}
Discuss separately the two cases in which the farthest-flying fragment after the lamp bursts does, and does not, strike the ceiling.
\end{hint}

\begin{hint}{1.05}
For the motion of the bullet, set up a three-dimensional coordinate system with the point of firing as origin; to find the locus of the aiming points, set up another two-dimensional coordinate system on the wall.
\end{hint}

\begin{hint}{1.06}
Note the velocity of the shell relative to the bomb, and discuss the relation between the minimum relative speed required for the shell to hit the bomb in the air and the speed with which the bomb is released.
\end{hint}

\begin{hint}{1.07}
The acceleration of $P$ relative to the paper is equal to its acceleration relative to the straight line.
\end{hint}

\begin{hint}{1.08}
The angular speed of the rod relative to the ground is equal to its angular speed relative to $A$ (that is, relative to the frame moving together with $A$).
\end{hint}

\begin{hint}{1.09}
The accelerations of hinges $B_3$ and $A_1$ are equal in magnitude, and the angular displacements of the rods are also equal.
\end{hint}

\begin{hint}{1.10}
First find the relative velocity of the centre of disc $B$ in the rotating frame of disc $A$; then use the relative-motion formula to find its absolute velocity.
\end{hint}

\begin{hint}{1.11}
In the rotating frame of the large steel ring, find the relative velocities and relative accelerations of points $A$ and $B$ on the small ring; the absolute accelerations of these two points then follow from the relative-motion formula.
\end{hint}

\begin{hint}{1.12}
Correctly finding the velocity and acceleration of the axis of the small cylinder relative to the ground, relative to cylinder $A$ and relative to cylinder $B$ is the key to this problem.
\end{hint}

\begin{hint}{1.13}
Draw the vector diagrams of the velocities and of the relative velocities of $A$ and $B$, and the diagram of the positions of $A$, $B$ and the centre $O$; then solve by geometry.
\end{hint}

\begin{hint}{1.14}
In the frame rotating with the rod, analyse whether the small ring has a position at which the resultant force on it vanishes; if it has, the ring need only reach this position in order to reach end $B$ of the rod.
\end{hint}

\begin{hint}{1.15}
The radius of curvature of the parabolic bent rod at a point can be found from the normal equation of a projectile motion whose trajectory has the same shape. Moreover, in the frame rotating with the rod, besides gravity, the normal reaction and the centrifugal force of inertia, the ring is also acted on by the Coriolis force because it is moving.
\end{hint}

\begin{hint}{1.16}
Analyse the motion, relative to the points of contact with the large disc, of two small segments on any diameter of the small disc that are symmetric about its centre; note that the friction forces on these two segments are directed opposite to this relative motion. Discuss the effect of the couple of these friction forces on the rotation of the small disc, and the angular speed of the small disc for which this frictional torque vanishes.
\end{hint}

\begin{hint}{1.17}
The three-ball system has two possible equilibrium configurations: the two lower balls separated, and the two lower balls close together.
\end{hint}

\begin{hint}{1.18}
When the large ring hangs on the nail in any position, the angle between the force exerted on it by the nail and the normal at that point is less than the angle of friction.
\end{hint}

\begin{hint}{1.19}
When $\alpha$ is the greatest value for which the system can be in equilibrium, there are two possible cases: in the first, ball 1 has already reached the critical state between sliding and not sliding relative to the inclined plane, while ball 2 has not yet reached this state; in the second, ball 2 has already reached this critical state, while ball 1 has not. Discuss the conditions for these two cases and the corresponding conclusions.
\end{hint}

\begin{hint}{1.20}
From the equilibrium conditions find the forces at the hinges at the ends of each rod, and hence the tension in the string; alternatively, use the principle of virtual work.
\end{hint}

\begin{hint}{1.21}
When the sphere is in the critical state between equilibrium and non-equilibrium, discuss which of the two contact points, that of the sphere with the board and that of the sphere with the inclined plane, has reached the critical state of relative sliding. Find the conditions for the two possible cases separately, and the corresponding value of $F$ follows.
\end{hint}

\begin{hint}{1.22}
After the small ball is fixed at the midpoint of rod $OA$, the symmetry of the system is destroyed. For the system to remain at rest, the minimum coefficient of friction sought must ensure not only that rod $OA$ does not slide relative to the ground, but also that the other two rods do not slide relative to the ground.
\end{hint}

\begin{hint}{1.23}
The direction of the normal reaction $N$ of the ring on the rod is perpendicular both to the rod and to the ring; the direction of the friction force $f$ acting on the rod is perpendicular both to the rod and to the direction of the normal reaction $N$ acting on it.
\end{hint}

\begin{hint}{1.24}
Note the continuity equation for the water, that is, the mass of water issuing from the fountain nozzle per second equals the mass of water rushing up to the bottom of the bucket.
\end{hint}

\begin{hint}{1.25}
By integrating twice the equation obtained from the differential form of the momentum theorem, find how the radius $r$ of the raindrop varies with time $t$; then find the velocity and the acceleration. The relations of $r$ and $t$ can also be found directly by solving the differential equation.
\end{hint}

\begin{hint}{1.26}
By combining the equation obtained from the differential form of the momentum theorem into an exact differential and then integrating, find the velocity.
\end{hint}

\begin{hint}{1.27}
After the ball leaves the table, change to the frame that moves upward with speed $v$ together with the upper end of the spring, and discuss the motion of the ball in this frame.
\end{hint}

\begin{hint}{1.28}
Note the motion of the centre of mass of the system; the force exerted by the light rod on ball $B$ can be found by analysing the motion of ball $A$ in the frame of ball $B$.
\end{hint}

\begin{hint}{1.29}
The speed of the ball is greatest when the resultant force along the tangent to its trajectory is zero.
\end{hint}

\begin{hint}{1.30}
When the ball slides to the right along the tube and ceases to interact with the tube, the ring has its greatest speed to the left; the radius of curvature of the trajectory at this point can be found by analysing the normal motion of the ball.
\end{hint}

\begin{hint}{1.31}
In the plane polar coordinate system with a certain bearing $a$ (or $b$) as pole, when the speed of body $c$ is greatest its acceleration is zero, hence its radial acceleration is also zero; also analyse the accelerations of bodies $A$ and $B$ at this instant.
\end{hint}

\begin{hint}{1.32}
In the centre-of-mass frame the two blocks vibrate while connected to a fixed point (the centre of mass) by an elastic rope; since the rope can only be stretched and not compressed, once it has returned to its natural length it exerts no force on the blocks.
\end{hint}

\begin{hint}{1.33}
Although after the supernova explosion the speed of the remaining star is unchanged, the centre of mass of the system formed by the remaining star and the other, unexploded star is no longer at rest. Moreover, in the centre-of-mass frame the mechanical energy of the system must be less than or equal to zero if the binary is not to separate.
\end{hint}

\begin{hint}{1.34}
Use conservation of momentum and the work-energy theorem.
\end{hint}

\begin{hint}{1.35}
Set up a plane polar coordinate system with the orthocentre of the equilateral triangle as pole, and in it discuss the boundary of the region of physically allowed points; when discussing the $\gamma$ photon, use the relativistic relation between the energy and momentum of a photon, namely $T=cP$.
\end{hint}

\begin{hint}{1.36}
The angular momentum of the ball about the vertical axis through the small hole (the symmetry axis of the conical surface) is conserved.
\end{hint}

\begin{hint}{1.37}
In the centre-of-mass frame the mechanical energy and the angular momentum of the system are both conserved; and when the spring reaches its maximum extension the ball has only a transverse velocity and no radial velocity.
\end{hint}

\begin{hint}{1.38}
The mechanical energy of a comet moving around the Sun on a parabolic orbit is zero; if the mechanical energy of the system is less than zero, the orbit is an ellipse or a circle.
\end{hint}

\begin{hint}{1.39}
The motion of the space shuttle (including the lunar module) around the Moon is similar to the motion of a planet around the Sun and can be solved by means of Kepler's laws.
\end{hint}

\begin{hint}{1.40}
The motion of the bullet can be analysed by means of a potential energy curve or an effective potential energy curve.
\end{hint}

\begin{hint}{1.41}
In the first scheme the motion of the spacecraft can be divided into two stages: the first stage is to escape the Earth's gravitational field, during which the Sun's gravitation may be regarded as unchanged; the second stage is then to escape from the solar system. In the second scheme the motion of the spacecraft can be divided into three stages: the first stage is the same as in the first scheme; the second stage is that, after escaping the Earth's gravitation, the spacecraft flies from the Earth's orbit to the orbit of Mars; the third stage is that, under the action of Mars's gravitation, it changes its direction of motion, leaves Mars's gravitational field tangentially from another point of the Martian orbit, and, with the aid of the orbital speed of Mars, continues to escape from the solar system.
\end{hint}

\begin{hint}{1.42}
After a planet receives a small radial perturbation it executes small radial oscillations; comparing the period of the small radial oscillation with the period of the original circular motion helps to describe the subsequent motion of the planet.
\end{hint}

\begin{hint}{1.43}
If the particle executes small radial oscillations, its orbit will be stable.
\end{hint}

\begin{hint}{1.44}
Whether the axis of rotation turns with constant angular velocity or is free to turn, the mechanical energy of the system in the rotating frame (including the centrifugal potential energy) is conserved; from conservation of mechanical energy the equation of oscillation of the ball is obtained.
\end{hint}

\begin{hint}{1.45}
The ball executes simple harmonic oscillations along the $x$ direction and along the $y$ direction with the same frequency but different amplitudes and different initial phases.
\end{hint}

\begin{hint}{1.46}
Compare which of the following two cases makes the lowest position reached by the cylinder lower: (1) before the cylinder descends to its lowest position the piston always remains inside the cylinder; (2) at some point before the cylinder descends to its lowest position the piston has already left the cylinder.
\end{hint}

\begin{hint}{1.47}
The angular momentum of the ball about the small hole $O$ is conserved; note that the rate of change of the radial velocity of the ball ($\ddot{r}$) is the acceleration of the up-and-down motion of the heavy object.
\end{hint}

\begin{hint}{1.48}
Find the tangential acceleration along the trajectory when the small ring moves along this steel wire, and compare this tangential acceleration with the acceleration of a particle executing simple harmonic motion along a straight line.
\end{hint}

\begin{hint}{1.49}
Because the wind speed has a gradient in the vertical direction perpendicular to the ground, when sound waves emitted from the ground travel between layers of different wind speeds, a refraction phenomenon similar to that occurring when light propagates in layers of different media appears, which changes their direction of propagation; after a long distance they bend back to the ground.
\end{hint}

\begin{hint}{1.50}
A wave reflected at a fixed end suffers a half-wave loss, that is, the oscillation excited by the reflected wave at the fixed end always has a phase difference of $\pi$ from that excited by the incident wave at the fixed end.
\end{hint}

\begin{hint}{1.51}
In the centre-of-mass frame the system has two normal modes: mode 1, $B$ is at rest while $A$ and $C$ move towards each other; mode 2, $A$ and $C$ move in the same direction while $B$ moves in the opposite direction.
\end{hint}

\begin{hint}{1.52}
This system has two normal modes: mode 1 is the symmetric mode, in which the rod remains horizontal and executes small oscillations along the vertical direction; mode 2 is the antisymmetric mode, in which the rod executes small angular oscillations about its centre of mass.
\end{hint}

\begin{hint}{1.53}
The trajectory of water ejected at any elevation angle $\theta$ is a quadratic equation in $u$ ($=\tan\theta$); for this equation to have a real solution its discriminant must be greater than or equal to zero, and hence the equation of the boundary curve of the region that the water jet can reach in space is obtained. The surface of revolution formed by rotating this curve about the axis of symmetry is the surface of the solid figure.
\end{hint}

\begin{hint}{1.54}
At depth $h$, the difference between the radial components of the gravitational force and of the centrifugal force of inertia acting on a small mass element is balanced by the force produced by the pressure difference between the two cross-sections along the radial direction. Moreover, if the centrifugal potential energy is introduced in the rotating frame, the total potential energy on the surface of the star is the same everywhere, and the value of $R_e-R_p$ is thereby easily found.
\end{hint}

\begin{hint}{1.55}
The surface of the Earth is an equipotential surface; this includes the Earth's gravitational potential, the Moon's gravitational potential, and the centrifugal potential in the centre-of-mass frame of the Moon--Earth system (the rotating frame).
\end{hint}

\begin{hint}{1.56}
From conservation of the angular momentum of the system formed by the planet and the satellite about the planet's centre, the relation between the variations of $\omega$, $\Omega$ and $r$ can be found; when the system reaches a steady state (a state in which the mechanical energy does not change), $\dfrac{\dd E}{\dd\omega}=0$ and $\dfrac{\dd^2E}{\dd\omega^2}>0$.
\end{hint}

\begin{hint}{1.57}
Bernoulli's equation gives the pressure at the points near the canopy, and the force on the canopy then follows from the pressure difference between the points inside and outside it.
\end{hint}

\begin{hint}{1.58}
At the instant of release, in the frame of the centre of mass of one rod, the rod turns about its centre of mass; relative to the ground the upper end of the rod moves vertically and its lower end moves horizontally; the various velocities then follow from the relative-motion relations.
\end{hint}

\begin{hint}{1.59}
Before the rod begins to slide relative to the corner of the table, the centre of mass of the rod moves in a circle; writing the tangential and normal equations of motion of the centre of mass gives the angle at which sliding begins.
\end{hint}

\begin{hint}{1.60}
In discussing the rotation of the cylinder, this cylinder differs from a solid cylinder of mass $m+M$, because the water inside does not take part in the rotation, so its moment of inertia is only that of the thin-walled cylinder. If the wedge-shaped block is not fixed, the force exerted by the cylinder on the block must be found by solving the motion of the cylinder in the frame of the block.
\end{hint}

\begin{hint}{1.61}
In the rotating frame the potential energy of the rod consists of the gravitational potential energy and the centrifugal potential energy; when the rod is in equilibrium the first derivative of the potential energy with respect to $\theta$ is zero, and the sign of the second derivative of the potential energy with respect to $\theta$ determines the stability of this equilibrium position.
\end{hint}

\begin{hint}{1.62}
The collision of a rotating sphere with the ground has two possible cases: in the first, the point of the sphere in contact with the ground keeps slipping during the whole collision (although the time is very short); in the second, relative slipping exists only during the first part of the collision and has ceased before the collision ends. In these two cases the impulse of the friction force is different, which leads to different horizontal velocities of the ball after the collision.
\end{hint}

\begin{hint}{1.63}
A rapidly spinning top has a large spin angular momentum; when its axis tilts, the torque about the centre of mass produced by the friction on the side of the point of contact between the axis and the ground makes the spin axis of the top gradually approach the vertical axis.
\end{hint}

\begin{hint}{1.64}
In the frame of the spacecraft the observer sees the Earth receding with speed $v$, so the speed of the radio signal emitted by the spacecraft relative to the Earth is $c-v$, and the time from the emission of the signal by the spacecraft to its reception on the Earth is equal to the time from the immediate retransmission by the Earth to the reception of this signal by the spacecraft.
\end{hint}

\begin{hint}{1.65}
The speed of light in the frame at rest relative to the glass is $\dfrac{c}{n}$; the speed of light in the moving glass in the ground frame can be found from the relativistic velocity-addition formula.
\end{hint}

\begin{hint}{1.66}
First find the $x$- and $y$-components of the velocity of the small packet thrown out in frame $S$; then, using the relativistic velocity transformation formula, find the $x'$- and $y'$-components of the velocity of the small packet in frame $S'$ of spacecraft $A$.
\end{hint}

\begin{hint}{1.67}
To the observer in frame $S'$, $B$ and $B'$ meet earlier than $A$ and $A'$; having found this time difference, one can find, from the condition that the reading of clock $B'$ is one unit of time when it receives the signal emitted by $A$, the relation between the proper length $l_0$ and $v$; while to the observer in frame $S$, $A$ and $A'$ meet earlier than $B$ and $B'$, and from the above relation between $l_0$ and $v$ the reading of clock $B$ when it receives the signal can be found.
\end{hint}

\begin{hint}{1.68}
The relativistic velocity transformation formula gives the angle between the direction from which light arrives and the direction of his flight, as seen by the pilot, for light coming from the edge of the Galaxy and from different parts of it, so as to determine the approximate shape of the Galaxy as he sees it; moreover, from the Doppler effect of light find the critical angle for redshift and blueshift, so as to determine who won the bet.
\end{hint}

\begin{hint}{1.69}
First find the proper wavelength of the light emitted by the source.
\end{hint}

\begin{hint}{1.70}
The decay of the particle is a process in which momentum and energy are conserved.
\end{hint}

\begin{hint}{1.71}
After finding the momentum and energy of the $\mu$ meson in the centre-of-mass frame, the momentum of the $\mu$ meson in the laboratory frame can be found by means of the relativistic momentum transformation formula.
\end{hint}

\begin{hint}{1.72}
The flight of the rocket from the Earth to the Galaxy is an accelerated process, so the rocket frame $S'$ is not an inertial frame. If at time $t$ a coordinate system is set up in $S'$ which moves uniformly relative to the Earth with the velocity $v$ of the rocket at time $t$ during an extremely short interval, then during this extremely short interval the frame $S'$ may be regarded as inertial, so the motion of the rocket can be discussed by means of special relativity, and the whole journey of the rocket is a combination of these infinitely short processes.
\end{hint}

\section{Thermodynamics}

\begin{hint}{2.01}
Use the equation of state of an ideal gas.
\end{hint}

\begin{hint}{2.02}
Use the equipartition theorem to find the mean potential energy of the molecular vibration.
\end{hint}

\begin{hint}{2.03}
Using the theoretical results of Brownian motion and the random-walk model, find the relation between the mean square value of the distance $x$ of a molecule from the origin and the time.
\end{hint}

\begin{hint}{2.04}
The number of collisions of molecules with unit area of the container per unit time is just the intensity of the molecular beam when the molecules leak out through the small hole; Clapeyron's equation.
\end{hint}

\begin{hint}{2.05}
For a van der Waals gas the internal energy is $U=C_VT-\dfrac{a}{V}+U_0$; the equation of an adiabatic process is $T(V-b)^{R/C_{V,m}}=$ const.
\end{hint}

\begin{hint}{2.06}
The polytropic process equation of one mole of a van der Waals gas: $\left(P+\dfrac{a}{V^2}\right)(V-b)^n=$ const.
\end{hint}

\begin{hint}{2.07}
Regard the atmosphere as an ideal gas and use the adiabatic process equation of an ideal gas to find the relation between the temperature $T$ and the height $z$.
\end{hint}

\begin{hint}{2.08}
The equation of state and the internal energy of a photon system are respectively $P=\dfrac13\sigma T^4$, $U=V\sigma T^4$.
\end{hint}

\begin{hint}{2.09}
Use Maxwell's distribution law to find the number of molecules colliding with the plate in time $\dd t$; note that the numbers of molecules colliding with the two faces of the plate are different.
\end{hint}

\begin{hint}{2.10}
Omitted.
\end{hint}

\begin{hint}{2.11}
$\lim_{T\to0}\dfrac{1}{T^2}\ee^{-1/T}=0$.
\end{hint}

\begin{hint}{2.12}
When the temperature range of variation is rather large, the latent heat of phase change $L$ depends on the temperature $T$.
\end{hint}

\begin{hint}{2.13}
Omitted.
\end{hint}

\begin{hint}{2.14}
This is a two-dimensional random-walk problem, usually discussed in mathematics by means of the multinomial distribution formula.
\end{hint}

\begin{hint}{2.15}
In a cavity at temperature $T_0$, a small black ball (regarded as a black body) at temperature $T$ transfers heat to the cavity wall at the rate $\dfrac{\dd Q}{\dd t}=A\sigma(T^4-T_0^4)$.
\end{hint}

\begin{hint}{2.16}
Omitted.
\end{hint}

\begin{hint}{2.17}
Omitted.
\end{hint}

\begin{hint}{2.18}
Omitted.
\end{hint}

\begin{hint}{2.19}
When the ideal gas in the container is in equilibrium, the mean translational kinetic energy of a molecule striking the wall is $\bar\varepsilon_{k\text{tr}}=2kT$.
\end{hint}

\begin{hint}{2.20}
When the volume of the liquid is constant, the work done by the surroundings on the liquid in an isothermal process is equal to the increase of the free energy of the liquid: $\dd W=\dd F=\sigma\,\dd A$.
\end{hint}

\begin{hint}{2.21}
At the interface where a liquid touches a solid, the contact angle made by the liquid surface is determined by the properties of the liquid and of the solid, and is independent of the position of the liquid surface.
\end{hint}

\begin{hint}{2.22}
The surface tension coefficient $\sigma$ of a given liquid is a function of the temperature $T$ only, and is independent of the area and the pressure of the liquid surface.
\end{hint}

\begin{hint}{2.23}
Omitted.
\end{hint}

\begin{hint}{2.24}
The relation between the entropy $S$ of a system and the thermodynamic probability $W^T$ is called the Boltzmann relation, or Boltzmann's hypothesis: $S=k\ln W^T$.
\end{hint}

\begin{hint}{2.25}
A one-dimensional chain means a chain that can be folded in one dimension.
\end{hint}

\begin{hint}{2.26}
Omitted.
\end{hint}

\begin{hint}{2.27}
Omitted.
\end{hint}

\begin{hint}{2.28}
Omitted.
\end{hint}

\section{Electromagnetism}

\begin{hint}{3.01}
(1) Consider the field of a point charge, taking a sphere centred on the point charge as a Gaussian surface. (2) Analyse the conservative nature of this electric force. (3) Consider the electric field at any point inside a uniformly charged spherical shell: take this point as vertex and construct a small cone, and examine the field produced at this point by the two surface elements of the shell cut out by the cone. (4) Take the conducting sphere as an example and use the result of (3).
\end{hint}

\begin{hint}{3.02}
(1) Use the energy method. (2) Regard the field of a quadrupole as the superposition of the fields of two oppositely directed dipoles, and consider the potential first.
\end{hint}

\begin{hint}{3.03}
(1) Relate the lines of force to the electric flux. (2) Examine the relation between the electric field (a vector) at a distant point and the two charges, noting that the distances from this point to the two charges are almost the same.
\end{hint}

\begin{hint}{3.04}
First take the approximation for a small quantity, then integrate.
\end{hint}

\begin{hint}{3.05}
(1) Omitted. (2) Divide the cube into eight equal small cubes; the centre of the original cube is then a vertex of the eight small cubes.
\end{hint}

\begin{hint}{3.06}
Use Gauss's law, taking as Gaussian surfaces concentric spheres of radii $r$ and $r+\dd r$ respectively.
\end{hint}

\begin{hint}{3.07}
Consider the force exerted by the positive ion layer on the negative ion layer and prove that this force is proportional to the small displacement. Replace the mass of the negative ion layer by the reduced mass to compensate for the motion of the positive ions.
\end{hint}

\begin{hint}{3.08}
(1) Note that the results of the integration inside and outside the spherical surface have different forms. (2) Since $\delta\ll1$, when the result of (1) is applied to the inner and outer spherical shells, take the small-quantity approximation.
\end{hint}

\begin{hint}{3.09}
Note the asymmetry of the charge on the two aluminium cylinders caused by accidental factors, and then analyse the positive feedback of this asymmetry.
\end{hint}

\begin{hint}{3.10}
To make the microphone give a voltage output, a condenser microphone must have a power supply, whereas an electret one needs none. In order that the motion of the diaphragm (which changes the capacitance) be reflected to the greatest extent in the output voltage, the charge on the diaphragm should be kept constant as far as possible.
\end{hint}

\begin{hint}{3.11}
(1) When a charge is moved, the image charge moves correspondingly, but the external force does work only on the charge being moved. (2) At this time the image charge no longer moves. (3) Consider the effects of (1) and (2) together.
\end{hint}

\begin{hint}{3.12}
Use the method of images; to find the position and magnitude of the image charge it suffices to consider special points on the axis of symmetry.
\end{hint}

\begin{hint}{3.13}
Regard these parallel conducting plates as capacitors connected in series.
\end{hint}

\begin{hint}{3.14}
(1) According to Ohm's law, write the differential equation for the relation between the radius and the time. (2) Note that external forces take part in doing work in the process.
\end{hint}

\begin{hint}{3.15}
Determine the charge on the third conducting sphere from the parallel connection of the capacitances, then determine the induced charge on the outer surface of the outer shell from electrostatic equilibrium; the third conductor is acted on by this charge.
\end{hint}

\begin{hint}{3.16}
First consider the induced charge distribution on an electrically neutral conducting sphere in a uniform external field (this can be obtained from two uniformly charged spheres of opposite sign displaced slightly in the external field, which makes the total field in the overlapping part vanish); the charge distribution on its upper half is the same as that on the conducting hemisphere in the problem, and the force on the hemisphere can then be found.
\end{hint}

\begin{hint}{3.17}
(1) At this time the charge may be regarded as uniformly distributed on the cylindrical surface, and the field in space is the field excited jointly by this charged cylindrical surface and its image. (2) Examine the equipotential surfaces of two parallel oppositely charged straight lines and prove that they are a family of non-concentric cylindrical surfaces; the problem is then solvable.
\end{hint}

\begin{hint}{3.18}
(1) Imagine the charge on the uniformly charged spherical surface to be frozen to the sphere. When the sphere is deformed along the $z$ axis into an ellipsoid of revolution, use the scaling principle to prove that the field inside is still everywhere zero, and hence obtain the charge distribution on the conducting ellipsoid of revolution. Find the potential at the centre by integration, and the capacitance follows. Note that when $b>a$ and when $b<a$ the results of the integration have different forms. (2) Find the small-quantity approximation separately.
\end{hint}

\begin{hint}{3.19}
(1) From the field inside the column, write the differential equation satisfied by the motion of the ions. (2) The charge $Q$ flowing to the conducting wire serves to cancel the reversed field formed by the ion layer.
\end{hint}

\begin{hint}{3.20}
Use the energy method, and regard the system as three spherical capacitors in series (the point charges being regarded as distributed on small conducting spheres), one of which is filled with a dielectric. When the point charge is moved to infinity, the three capacitors may still be regarded as in series, but all are then vacuum capacitors.
\end{hint}

\begin{hint}{3.21}
Suppose an electric field exists between the plates; find the polarisation intensity of each dielectric sphere under the action of the external field (neglecting the influence of the other dielectric spheres), and hence the electric moment of each dielectric particle, and then the equivalent polarisation intensity and the equivalent permittivity of the dielectric made of sparse particles.
\end{hint}

\begin{hint}{3.22}
First, since polarisation charges of opposite sign gather around the point charge, the behaviour of the point charge $q$ is equivalent to that of a point charge of magnitude $q/\varepsilon_r$. Then use the boundary conditions at the surface of the dielectric to find the distribution of the surface polarisation charge; the behaviour of the surface polarisation charge is equivalent to that of an ``image charge''.
\end{hint}

\begin{hint}{3.23}
(1) Use Gauss's law. (2) Use the method of the solution to Problem 3.21 to find the polarisation intensity and the dipole moment of the spherical dust particle, and hence the force on the dust. (3) Write the equation of motion of the dust and integrate.
\end{hint}

\begin{hint}{3.24}
(1) The spherical cavity is a uniform field, which can be found from the boundary conditions. (2) Use the boundary conditions to find the charge distribution on the surface of the conducting sphere and on the surface of the dielectric cavity, and use the energy method to find the force on the conducting sphere.
\end{hint}

\begin{hint}{3.25}
The cavity in which $q'$ is located lies in the field excited by $q$ and by the polarisation charge on the surface of the cavity in which it is located (this field may be regarded as uniform at the cavity of $q'$); use the result of part (1) of Problem 3.24.
\end{hint}

\begin{hint}{3.26}
The distribution of the polarisation charge on the surface of the dielectric sphere is the same as that of a uniformly polarised dielectric sphere.
\end{hint}

\begin{hint}{3.27}
The current distribution in the resistive material may be regarded as the superposition of the current distribution when only the current flows in and that when only the current flows out. From the differential form of Ohm's law it follows that the electric field in the material is equivalent to the field of a pair of point charges of opposite sign.
\end{hint}

\begin{hint}{3.28}
Using similarity and scaling relations, the equipotential points and the superposition principle, simplify the original circuit to the simple circuit shown in the figure.
\end{hint}

\begin{hint}{3.29}
This problem can be solved by Kirchhoff's method, by the superposition of currents, or by the star--delta transformation.
\end{hint}

\begin{hint}{3.30}
Use Thévenin's theorem and the small-quantity approximation.
\end{hint}

\begin{hint}{3.31}
(1) Use Thévenin's theorem repeatedly from left to right. (2) Looking to the left from $A_i$ gives an equivalent source, and looking to the right from $A_{i-1}$ also gives an equivalent source; the original circuit is thus simplified to two equivalent sources and the resistance $R$ between $A_i$ and $A_{i-1}$ connected in series.
\end{hint}

\begin{hint}{3.32}
(1) This can be decided by whether the directions of the current in the ring and of the velocity of the ring are the same. (2) When the ring decelerates in its rotation, the current carriers are acted on by an inertial force. This inertial force may be regarded as a force of non-electrostatic origin, from which an equivalent electromotive force is obtained; using Ohm's law and integrating the current over time, the specific charge can be found.
\end{hint}

\begin{hint}{3.33}
(1) Charge will accumulate on the surface of the bubble; the surface charge distribution can be found from the boundary conditions for the current and for the electric field. The field inside the bubble is the superposition of the original field in the conductor and the field produced by the surface charge (a uniform field). (2) The field at any point $P$ inside the conductor is the superposition of the original field and the field produced by the surface charge (the field of a dipole).
\end{hint}

\begin{hint}{3.34}
(1) Regard the current as steady; the current density depends only on the radius, and the electric field follows from Ohm's law. (2) Use Gauss's law. (3) The line integral of the electric field along a path gives the potential of the Earth's surface. (4) Write the differential equation for the relation between the charge and the time.
\end{hint}

\begin{hint}{3.35}
(1) Use Gauss's law and Ampère's circuital law to find the electric and magnetic fields; from the statement of the problem (the electric and magnetic forces cancel) find the power. (2) The self-inductance $L^*$ per unit length of the transmission line can be found from the capacitance $C^*$ per unit length ($L^*C^*=\varepsilon_0\mu_0$).
\end{hint}

\begin{hint}{3.36}
(1) The field can be obtained by superposing (reducing to an integral) the magnetic fields produced on the axis by circular coils. (2) It can be found from the variation of the magnetic field along the axis by using Gauss's law for the magnetic field.
\end{hint}

\begin{hint}{3.37}
The dielectric is polarised under the joint action of the Lorentz force and the electric field produced by its own polarisation charges.
\end{hint}

\begin{hint}{3.38}
Same as Hint 3.37. Use Gauss's law to find the charge distribution.
\end{hint}

\begin{hint}{3.39}
(2) In producing a current the Lorentz force is equivalent to the electric force; the current intensity between the inner and outer cylinders is determined by the radial component of the current density, while the radial component of the Lorentz force is determined by the transverse component of the drift velocity of the carriers, and the transverse component of the Lorentz force is in turn determined by the radial component of the drift velocity; a relation between the current density, the radial momentum and the electric and magnetic fields can thus be established.
\end{hint}

\begin{hint}{3.40}
(1) The dielectric sphere will be uniformly magnetised in the magnetic field at the centre of the current-carrying coil. (2) The magnetised dielectric sphere may be regarded as a magnetic dipole; it will increase the magnetic flux through the coil.
\end{hint}

\begin{hint}{3.41}
(1) The motion of charge produces a current. (2) A changing magnetic field produces an electric field.
\end{hint}

\begin{hint}{3.42}
The long rod will be uniformly magnetised along the $z$ direction; the magnetisation current in the rod produces a uniform magnetic field inside the rod, and outside the rod it is equivalent to the field produced by two opposite currents.
\end{hint}

\begin{hint}{3.43}
Use the angular momentum theorem and conservation of energy. The torque acting on the electron comes from the Lorentz force.
\end{hint}

\begin{hint}{3.44}
(1) The centripetal force for the circular motion of the particle is the Lorentz force. (2) Since the magnetic flux density decreases along the radial direction, the lines of magnetic induction must be barrel-shaped; this makes the magnetic flux density have a radial component at $z>0$ and at $z<0$, so that when the particle is slightly displaced in the vertical direction it experiences a restoring force.
\end{hint}

\begin{hint}{3.45}
The particle moves in a helix about the $z$ axis; since the magnetic field gradually increases along the $z$ direction, a radial component of the magnetic field appears. The Lorentz force produced by this radial component points in the $-z$ direction, so that the particle is finally reflected. Using Newton's law and conservation of energy, establish the equation relating the rate of change of $v_z$ with $z$ to the rate of change of $B_z$ with $z$.
\end{hint}

\begin{hint}{3.46}
In the two cases (1) and (2), write the equation of motion of the coil and the circuit equation separately.
\end{hint}

\begin{hint}{3.47}
(1) Take the time of falling over as a parameter and find the ratio of the kinetic energy to the Joule heat. The time of falling over can be estimated from the case without a magnetic field. (2) The time taken for the ring to fall over is found by balancing the moment of the gravitational force against the moment of the Ampere force.
\end{hint}

\begin{hint}{3.48}
(1) An induced current will appear in the superconducting coil so as to keep the magnetic flux through the coil zero. This current, acting in the radial magnetic field, will be acted on by an upward vertical force that balances gravity, and the radial component of the magnetic field can be found from the rate of change of the $z$ component of the magnetic field along the $z$ direction. (2) From the condition that the magnetic flux through the coil is zero, find the variation of the induced current in the coil with position, then write the equation of motion of the coil; note the small-quantity approximation.
\end{hint}

\begin{hint}{3.49}
Reduce the steel magnet to a magnetic dipole and the metal tube to a series of coaxial circular coils; use conservation of energy.
\end{hint}

\begin{hint}{3.50}
Use the phasor method.
\end{hint}

\begin{hint}{3.51}
Write Kirchhoff's equations (using the complex method).
\end{hint}

\begin{hint}{3.52}
Use the complex method together with the boundary condition at infinity.
\end{hint}

\begin{hint}{3.53}
Use the quasi-stationary condition, and note that the displacement current is the same as the conduction current in exciting a magnetic field.
\end{hint}

\begin{hint}{3.54}
(2) When the pulse is applied, the capacitor may be regarded as a short circuit; it is charged by the current $U_0/R_1$ and then discharges through the parallel combination of $R_1$ and $R_2$.
\end{hint}

\begin{hint}{3.55}
(1) Adjust $R$ to the critically damped state. (2) When the current is maximum, $\dd I/\dd t=0$; use the final-state condition.
\end{hint}

\begin{hint}{3.56}
(1) Express the propagation factor in complex form, then use the condition $|a|<1$ to find $\omega_c$. (2) Use the boundary condition at infinity. (3) Examine the propagation factor between two adjacent sections, $a_n=U_{n+1}/U_n$ ($n$ the section number), and set $|a_n|=1$. (4) Under the condition $\omega\ll\omega_c$, find the phase difference $\delta$ between adjacent sections; then the speed of propagation of the signal along the network is $v=\Delta r/(\delta/\omega)$.
\end{hint}

\begin{hint}{3.57}
Use the star--delta transformation of complex impedances.
\end{hint}

\begin{hint}{3.58}
(1) Superpose the incident and reflected waves, and use $E=0$ inside an ideal conductor and the boundary conditions for the electric and magnetic fields. (2) Apply Ampère's circuital law at the surface of the conductor to find the surface current, then use Ampère's law to find the force on the surface of the conductor.
\end{hint}

\section{Optics}

\begin{hint}{4.01}
(1) Express the deviation angle $\delta$ as $\delta=i+r'-A$ ($r'$ the angle of emergence), and use the law of refraction to express $\delta$ as a function of the angle of incidence $i$. (2) Set $\dd\delta/\dd i=0$; the result then follows.
\end{hint}

\begin{hint}{4.02}
(1) Use the symmetry of the optical path and the law of refraction. (2) Calculate successively with a calculator.
\end{hint}

\begin{hint}{4.03}
Take the $xy$ plane as the cross-section and the $x$ axis as the axis of rotation; using the equality of the optical paths of an object and its image, write the equation of the cross-section curve, then simplify and discuss.
\end{hint}

\begin{hint}{4.04}
First use the law of refraction and geometrical relations to express the complement $\alpha$ of the deviation angle $\delta$ as a function of the angle of incidence $i$ and the angle of refraction $r$ (see the figure); $\dd\alpha/\dd i=0$ determines $\alpha$, because in that direction the light is most concentrated; $\alpha$ is the angular radius of the rainbow.
\end{hint}

\begin{hint}{4.05}
Use the law of refraction and the small-angle approximation ($\alpha\ll1$, the ray is bent only slightly), then integrate.
\end{hint}

\begin{hint}{4.06}
(1) Use the law of refraction and the small-angle approximation [($\alpha-\alpha_0$) is a small angle]. (2) Omitted.
\end{hint}

\begin{hint}{4.07}
Regard the laser beam as a stream of photons; the change of momentum caused by the deflection of the photons by the lens immediately produces a force on the lens.
\end{hint}

\begin{hint}{4.08}
Regard the thin convex lens as a combination of two spherical surfaces.
\end{hint}

\begin{hint}{4.09}
The diameter of the light spot on screen $Q$ is determined by the position of the image point and by the angular aperture of the emergent beam, and the angular aperture of the emergent beam is determined by the size of the hole in screen $P$; more precisely, it is determined by the size of the image of the hole formed by lens $L_2$.
\end{hint}

\begin{hint}{4.10}
(1) From the object--image correspondence, when the concave mirror is at the position of the image of object $O$ formed by the lens, the final image will coincide with the original object. Another position of the concave mirror can be found from the requirement that a ray emitted from $O$ parallel to the optical axis retraces its path after passing through the lens. Of course, the two positions of the concave mirror can also be found by writing the equation from the condition that the magnification is $+1$. (2) Omitted.
\end{hint}

\begin{hint}{4.11}
(1) Use the law of refraction and the sine rule for triangles; alternatively, first find the position of the image by paraxial rays. (2) 1) From the position of the aplanatic point obtained in part (1) and the conditions given in the problem, $R_2$ and $R_3$ can be found (the second lens being concave--convex). 2), 3), 4) Omitted.
\end{hint}

\begin{hint}{4.12}
(1) Use the small-angle approximation to find the image of $S$ formed by the upper and lower prisms (note that the angles of incidence and refraction of the ray travelling from $S$ to the prisms are also small). (2) Find the optical path difference of the two rays passing through the upper and lower prisms after $S$ has been displaced upward.
\end{hint}

\begin{hint}{4.13}
(1) Note the phase jump that occurs when light passes through a focal point. (2) Omitted.
\end{hint}

\begin{hint}{4.14}
(1) Use the reversibility of time (that is, the reversibility of the optical path) and the superposition principle. (2) Superpose the complex amplitudes of the reflected waves and use the result of part (1).
\end{hint}

\begin{hint}{4.15}
Use the relativistic velocity-addition formula.
\end{hint}

\begin{hint}{4.16}
Sum separately the complex amplitudes of the terms with coefficient $1$ and with coefficient $\alpha$ in the reflectivity $(1\pm\alpha)$.
\end{hint}

\begin{hint}{4.17}
Regard each step of the sailor as a vector; the problem reduces to finding the magnitude of the sum (resultant) of a large number of vectors of equal magnitude whose directions vary at random.
\end{hint}

\begin{hint}{4.18}
Regard the light disturbance at any point near the geometrical image point on the observation screen as the superposition (integral) of the light disturbances produced at that point by all the surface elements of the geometrical wavefront formed by the square aperture.
\end{hint}

\begin{hint}{4.19}
(1) The incident light may be regarded as the superposition of unpolarised light and elliptically polarised light; the $x$ and $y$ axes are then the principal axes of the ellipse, and the phase difference between the components of the vibration of the elliptically polarised light along these two axes is $90\degree$. (2) Elliptically polarised light becomes linearly polarised after passing through a quarter-wave plate.
\end{hint}

\begin{hint}{4.20}
In a birefringent crystal, light polarised parallel and perpendicular to the optic axis travels with different speeds; after passing through a crystal of a certain thickness, these two polarised components acquire a phase difference proportional to the thickness of the crystal, and when they pass through the analyser behind, they again combine into linearly polarised light.
\end{hint}

\section{Atomic Physics}

\begin{hint}{5.01}
Use Einstein's photoelectric equation.
\end{hint}

\begin{hint}{5.02}
If the particle is a $\gamma$ photon, the phenomenon in (b) is the Compton effect of the proton; comparing the experimental results of (a) and (b), one sees that only by abandoning conservation of energy and momentum can the hypothesis that the particle is a $\gamma$ ray be supported.
\end{hint}

\begin{hint}{5.03}
The Compton effect and the inverse Compton effect.
\end{hint}

\begin{hint}{5.04}
Conservation of momentum and energy, and the relativistic expressions for the momentum--energy relation.
\end{hint}

\begin{hint}{5.05}
Estimate by means of the uncertainty relation. Since the kinetic energy of the electron is comparable to its rest energy, the relativistic momentum--energy relation should be used for the electron.
\end{hint}

\begin{hint}{5.06}
Use the uncertainty relation $\Delta x\,\Delta p\geq\hbar$ to eliminate $r$ or $p$ from the energy expression.
\end{hint}

\begin{hint}{5.07}
Use the quantisation condition for the orbital angular momentum of the electron.
\end{hint}

\begin{hint}{5.08}
Hydrogen-like ion spectra; take the finite nuclear mass correction into account.
\end{hint}

\begin{hint}{5.09}
Take the finite nuclear mass correction into account.
\end{hint}

\begin{hint}{5.10}
The influence of the finite size of the nucleus on the energy levels and spectrum of the exotic atom.
\end{hint}

\begin{hint}{5.11}
The atomic core of the sodium atom may be regarded as an effective nucleus; the valence electron is subject to different effective nuclear charges in different states.
\end{hint}

\begin{hint}{5.12}
The answers to (1) and (2) are helpful for the solution of (3).
\end{hint}

\begin{hint}{5.13}
The Franck--Hertz experiment.
\end{hint}

\begin{hint}{5.14}
Use the Bohr model of the hydrogen atom to estimate the energies of some states of the helium atom (a two-electron atom).
\end{hint}

\begin{hint}{5.15}
This problem mainly uses the fact that the uranium series reaches secular equilibrium after a certain time, that is, the activities of the parent and of all the daughters are equal, and the assumption that the age of formation of the uranium ore is approximately the age of the Earth.
\end{hint}

\begin{hint}{5.16}
The age of ancient biological samples can be determined from the exponential law of variation of activity with time.
\end{hint}

\begin{hint}{5.17}
The key is to write down the energy requirement for the excited state $^{17}\mathrm{O}^*$ to emit a neutron and reach $^{16}\mathrm{O}$ (which may be in the ground state or in an excited state).
\end{hint}

\begin{hint}{5.18}
Using the Coulomb energy formula directly related to the same mass number $A$ and radius $R$, $B_c=\dfrac35\dfrac{z(z-1)e^2}{4\pi\varepsilon_0R}$, the corresponding nuclear radius constant can be found (note that the nuclear radius constants for light and heavy nuclei actually differ somewhat).
\end{hint}

\begin{hint}{5.19}
Use the characteristics of elastic collisions in the centre-of-mass frame and the conversion relation between the differential cross-section in the centre-of-mass frame and that in the laboratory frame.
\end{hint}
